%%═══════════════════════════════════════════════════════════════════
%% Lecture 05 — Radioactive Decay, Activity, and Radiation Dose
%% 77603 — Nuclear Physics
%% Date: 12 May 2026
%% Lecturer: Dr. Moshe Friedman
%%═══════════════════════════════════════════════════════════════════

\section{Lecture 05 — Radioactive Decay, Activity, and Radiation Dose}
\label{sec:lec05}
\date{12 May 2026}

%%───────────────────────────────────────────────────────────────────
\subsection*{Overview}
%%───────────────────────────────────────────────────────────────────
This lecture opens the second part of the course.  The first part treated the
nucleus as a bound quantum system: nuclear sizes, binding energies, scattering,
the deuteron, and shell structure.  We now turn to the processes by which an
unstable nucleus changes in time.  The central idea is simple but powerful:
radioactive decay is a \emph{statistical quantum process} in which a nucleus
loses energy by moving to a more favourable nuclear state, often by becoming a
different nucleus.

The lecture has four main goals.  First, we classify the most important decay
modes: $\gamma$ decay, internal conversion, $\alpha$ decay, fission, $\beta$
decay, electron capture, and delayed particle emission.  Second, we quantify
radioactivity using the decay constant, activity, half-life, and mean lifetime.
Third, we study decay branches and decay chains, where a radioactive parent
produces radioactive daughters.  Finally, we introduce the language used in
radiation safety: exposure, absorbed dose, equivalent dose, and the practical
principles ALARA and LNT.

%%───────────────────────────────────────────────────────────────────
\subsection{What Counts as Radioactive Decay?}
\label{subsec:lec05:definition}

\subsubsection*{Physical Motivation}
A nucleus can be energetically unstable even when it is a perfectly well-defined
bound quantum state.  It may sit in a local minimum: classically it looks stable,
but quantum mechanically there can be a small probability per unit time to
escape to a lower-energy configuration.  This is why a nucleus can have a
lifetime ranging from less than a nanosecond to longer than the age of the
Universe.

\dfn{Radioactive decay}{Radioactive decay is a process in which a nucleus in its ground state or in an
excited state releases energy and moves to a lower-energy nuclear configuration.
The final configuration may be the same nucleus in a lower internal state, or a
different nucleus.}

The energy scale is set by nuclear binding energies.  In most ordinary nuclear
decays the released energy is of order
\begin{equation}
  \label{eq:lec05:decay_energy_scale}
  \boxed{Q \sim 10^2\,\mathrm{keV}\text{--}10\,\mathrm{MeV}.}
\end{equation}
\textit{Physical significance:} This is much larger than atomic transition
energies, but much smaller than the rest energy of a nucleus.  Fission is the
main exception: a single fission event releases roughly $200\,\mathrm{MeV}$.

Because decay is probabilistic, we should not ask, ``When will this particular
nucleus decay?'' as if the answer were deterministic.  The meaningful question
is: \emph{what is the probability per unit time that it decays?}  That question
leads directly to the decay constant $\lambda$ in
Section~\ref{subsec:lec05:decay_law}.

%%───────────────────────────────────────────────────────────────────
\subsection{Decay Modes}
\label{subsec:lec05:decay_modes}

\subsubsection*{Physical Motivation}
The nucleus must get rid of energy while preserving the conservation laws:
energy, momentum, angular momentum, electric charge, baryon number, and lepton
number where relevant.  Different decay modes are different ways of satisfying
those constraints.  Some modes keep the identity of the nucleus fixed; others
change $Z$, $A$, or both.

\paragraph{Gamma decay.}
In $\gamma$ decay an excited nucleus emits a high-energy photon and remains the
same nuclide:
\begin{equation}
  \label{eq:lec05:gamma_decay}
  \boxed{{}^{A}_{Z}X^{*}\rightarrow{}^{A}_{Z}X+\gamma.}
\end{equation}
The photon carries almost all the transition energy, but not exactly all of it:
the daughter nucleus must recoil so that momentum is conserved.

\textit{Physical significance:} A laboratory ``gamma source'' is usually not a
source whose nuclei sit for macroscopic times in excited states.  Excited
nuclear states typically live for very short times, often nanoseconds or less.
The gamma photon is usually produced after some earlier decay, such as
$\beta$ decay, leaves the daughter nucleus excited.

\paragraph{Internal conversion.}
An excited nucleus can also give its energy directly to an atomic electron,
which is then ejected:
\begin{equation}
  \label{eq:lec05:internal_conversion}
  \boxed{{}^{A}_{Z}X^{*}\rightarrow{}^{A}_{Z}X^{+}+e^{-}.}
\end{equation}
The ion charge reminds us that an atomic electron has been removed.  In practice
one often focuses on the nuclear transition and treats the atomic rearrangement
separately.

\textit{Physical significance:} Internal conversion competes with gamma
emission because both start from the same excited nuclear state.  Which channel
dominates depends on the transition energy, multipolarity, and the electron
wavefunction near the nucleus.

\paragraph{Alpha decay.}
In $\alpha$ decay a heavy nucleus emits a helium-4 nucleus:
\begin{equation}
  \label{eq:lec05:alpha_decay}
  \boxed{{}^{A}_{Z}X\rightarrow{}^{A-4}_{Z-2}Y+{}^{4}_{2}\mathrm{He}.}
\end{equation}
The mass number decreases by $4$ and the proton number decreases by $2$.
Although the $\alpha$ particle is strongly bound, it must typically tunnel
through the Coulomb barrier to escape the parent nucleus.

\textit{Physical significance:} Alpha decay is a clean example of a nucleus
being energetically allowed to move to a lower state, while still being
classically trapped for a long time.  The lifetime is extremely sensitive to the
barrier width and to the decay energy.

\paragraph{Fission.}
In fission a very heavy nucleus splits into two medium-mass fragments, usually
with extra neutrons and gamma rays:
\begin{equation}
  \label{eq:lec05:fission_decay}
  \boxed{{}^{A}_{Z}X\rightarrow{}^{A_1}_{Z_1}Y
  +{}^{A_2}_{Z_2}W+\bar{\nu}\,n+\gamma+\text{kinetic energy}.}
\end{equation}
Here $\bar{\nu}$ is the average number of emitted neutrons.  The notation is
schematic: real fission has many possible fragment pairs and neutron
multiplicities.

\textit{Physical significance:} Fission releases an unusually large energy,
about $200\,\mathrm{MeV}$ per event, because the fragments lie closer to the
maximum of binding energy per nucleon than the original very heavy nucleus.

\paragraph{Beta-minus decay.}
In $\beta^{-}$ decay, a neutron inside the nucleus turns into a proton:
\begin{equation}
  \label{eq:lec05:beta_minus_decay}
  \boxed{{}^{A}_{Z}X\rightarrow{}^{A}_{Z+1}Y+e^{-}+\bar{\nu}_{e}.}
\end{equation}
The emitted electron preserves electric charge, and the electron antineutrino
preserves lepton number.  At the nucleon level,
\begin{equation}
  \label{eq:lec05:neutron_beta_decay}
  n\rightarrow p+e^{-}+\bar{\nu}_{e}.
\end{equation}

\textit{Physical significance:} Because the final state contains three bodies
(the daughter nucleus, electron, and antineutrino), the electron energy is not a
single sharp value.  It forms a continuous spectrum from near zero up to the
endpoint energy.

\paragraph{Beta-plus decay.}
In $\beta^{+}$ decay, a proton inside the nucleus turns into a neutron:
\begin{equation}
  \label{eq:lec05:beta_plus_decay}
  \boxed{{}^{A}_{Z}X\rightarrow{}^{A}_{Z-1}Y+e^{+}+\nu_{e}.}
\end{equation}
The positron will usually annihilate with an electron, producing two photons:
\begin{equation}
  \label{eq:lec05:positron_annihilation}
  e^{+}+e^{-}\rightarrow\gamma+\gamma,
  \qquad E_{\gamma}=511\,\mathrm{keV}.
\end{equation}
The two photons are emitted nearly back-to-back in the centre-of-mass frame.

\textit{Physical significance:} This back-to-back photon signature is the basis
of positron emission tomography (PET): a positron-emitting isotope is placed in
the body, and detectors reconstruct where the annihilations occurred.

\paragraph{Electron capture.}
Electron capture competes with $\beta^{+}$ decay.  Instead of emitting a
positron, the nucleus captures one of its atomic electrons:
\begin{equation}
  \label{eq:lec05:electron_capture}
  \boxed{{}^{A}_{Z}X+e^{-}\rightarrow{}^{A}_{Z-1}Y+\nu_{e}.}
\end{equation}
At the nucleon level, a proton and electron combine into a neutron and a
neutrino.

\textit{Physical significance:} Electron capture and $\beta^{+}$ decay both
reduce $Z$ by one while keeping $A$ fixed.  They are therefore competing
channels for proton-rich nuclei.

\paragraph{Delayed proton and neutron emission.}
Some nuclei are produced in excited states above the particle-emission
threshold.  They can then emit a proton or neutron:
\begin{align}
  \label{eq:lec05:proton_emission}
  {}^{A}_{Z}X^{*} &\rightarrow{}^{A-1}_{Z-1}Y+p,\\
  \label{eq:lec05:neutron_emission}
  {}^{A}_{Z}X^{*} &\rightarrow{}^{A-1}_{Z}Y+n.
\end{align}
These decays are less common in naturally occurring radioactivity, but they are
important near the proton and neutron drip lines and in accelerator-produced
nuclei.

\textit{Physical significance:} Particle emission reminds us that ``decay
mode'' is not only about the parent nucleus.  It also depends on which daughter
states are energetically accessible.

%%───────────────────────────────────────────────────────────────────
\subsection{Branches, Cascades, and Why Gamma Rays Often Come Last}
\label{subsec:lec05:branches}

\subsubsection*{Physical Motivation}
A nucleus does not necessarily have only one allowed decay path.  The same
parent state can decay to several daughter states, and a daughter state may
itself be unstable.  Nuclear data tables therefore show arrows with
percentages: these are branching ratios.

If a parent can decay through two channels with partial decay constants
\(\lambda_a\) and \(\lambda_b\), the total decay constant is
\begin{equation}
  \label{eq:lec05:partial_lambdas}
  \boxed{\lambda=\lambda_a+\lambda_b.}
\end{equation}
The branching ratios are
\begin{equation}
  \label{eq:lec05:branching_ratios}
  \boxed{b_a=\frac{\lambda_a}{\lambda},\qquad
  b_b=\frac{\lambda_b}{\lambda},\qquad b_a+b_b=1.}
\end{equation}
\textit{Physical significance:} The nucleus decays with the total rate
\(\lambda\), but once a decay occurs the probability of a particular channel is
set by its partial rate divided by the total rate.

If the daughter is created in an excited state, it may emit one or more gamma
photons while stepping down through the level scheme.  Such a sequence is called
a cascade.  This is why gamma rays often appear as secondary radiation after a
beta or alpha decay.

%%───────────────────────────────────────────────────────────────────
\subsection{The Radioactive Decay Law}
\label{subsec:lec05:decay_law}

\subsubsection*{Physical Motivation}
The defining assumption is that each nucleus has the same probability per unit
time to decay, independent of how long it has already existed.  This is a
memoryless process.  The macroscopic exponential law follows from that one
microscopic statement.

Let \(N(t)\) be the number of undecayed nuclei at time \(t\).  If each nucleus
has probability \(\lambda\,dt\) to decay during the short interval \(dt\), then
\begin{equation}
  \label{eq:lec05:decay_differential}
  \boxed{\dv{N}{t}=-\lambda N.}
\end{equation}
Solving the differential equation gives
\begin{equation}
  \label{eq:lec05:decay_solution}
  \boxed{N(t)=N_0 e^{-\lambda t}.}
\end{equation}
\textit{Physical significance:} The decay curve is exponential because the
loss rate is proportional to how many unstable nuclei remain.

The activity is the number of decays per unit time:
\begin{equation}
  \label{eq:lec05:activity_definition}
  \boxed{\activity(t)=\lambda N(t)=\lambda N_0e^{-\lambda t}.}
\end{equation}
Activity has units of inverse time:
\begin{equation}
  \label{eq:lec05:activity_units}
  \boxed{1\,\mathrm{Bq}=1\,\mathrm{s}^{-1}=1\,\mathrm{decay/s},\qquad
  1\,\mathrm{Ci}=3.7\times10^{10}\,\mathrm{Bq}.}
\end{equation}
\textit{Physical significance:} Activity measures how many decays occur per
second, not how much biological damage is caused.  A high-activity source is
not fully characterized until we also know the radiation type and energy.

\subsubsection{Half-Life and Mean Lifetime}
\label{subsubsec:lec05:half_life}

The half-life \(T_{1/2}\) is defined by \(N(T_{1/2})=N_0/2\).  Substituting
Eq.~\eqref{eq:lec05:decay_solution},
\begin{equation}
  \label{eq:lec05:half_life_derivation}
  \frac{1}{2}=e^{-\lambda T_{1/2}}
  \quad\Rightarrow\quad
  \boxed{T_{1/2}=\frac{\ln2}{\lambda}.}
\end{equation}
The mean lifetime \(\tau\) is the average decay time of one nucleus:
\begin{equation}
  \label{eq:lec05:mean_lifetime}
  \boxed{\tau=\frac{\int_0^\infty t\,\lambda e^{-\lambda t}\,dt}
  {\int_0^\infty \lambda e^{-\lambda t}\,dt}
  =\frac{1}{\lambda}.}
\end{equation}
Thus
\begin{equation}
  \label{eq:lec05:tau_half_life_relation}
  \boxed{T_{1/2}=(\ln2)\tau.}
\end{equation}
\textit{Physical significance:} ``Lifetime'' is sometimes used loosely in
speech.  In formulas, \(T_{1/2}\) and \(\tau\) are different by the factor
\(\ln2\).

%%───────────────────────────────────────────────────────────────────
\subsection{How Decay Constants Are Measured}
\label{subsec:lec05:measuring_lifetimes}

\subsubsection*{Physical Motivation}
A lifetime measurement is not just a curve drawn through points.  Radioactive
decay is intrinsically statistical, so the scatter of the data is part of the
physics.  A good fit should not have every point exactly on the curve; it should
have deviations consistent with the measurement uncertainties.

In a simple lifetime measurement, one produces a radioactive isotope, records
the number of decays in time bins, and fits the counting rate to
\begin{equation}
  \label{eq:lec05:counting_fit}
  \boxed{\activity(t)=\activity_0e^{-\lambda t}+B,}
\end{equation}
where \(B\) is a possible background rate.  The quality of the fit is often
checked with
\begin{equation}
  \label{eq:lec05:chi_squared}
  \boxed{\chi^2=\sum_i
  \left(\frac{y_i-f(x_i)}{\Delta y_i}\right)^2.}
\end{equation}
If the uncertainties \(\Delta y_i\) are realistic and the model is correct,
\(\chi^2\) should be of order the number of degrees of freedom
\(\nu=N_\mathrm{points}-N_\mathrm{fit\,parameters}\).

\textit{Physical significance:} The usual correlation coefficient \(R^2\) is
not enough for a physics fit because it does not know the experimental
uncertainties.  A meaningful goodness-of-fit test must compare residuals to
their errors.

%%───────────────────────────────────────────────────────────────────
\subsection{Building Radioactivity: Activation}
\label{subsec:lec05:activation}

\subsubsection*{Physical Motivation}
Radioactivity can be produced artificially.  A stable material can be irradiated
by particles or photons so that some nuclei are converted into radioactive
isotopes.  This is how many medical radioisotopes are produced in reactors or
accelerators.

Suppose a radioactive isotope is produced at a constant rate \(R\) nuclei per
second and decays with constant \(\lambda\).  Then
\begin{equation}
  \label{eq:lec05:activation_differential}
  \dv{N}{t}=R-\lambda N,\qquad N(0)=0.
\end{equation}
Solving,
\begin{equation}
  \label{eq:lec05:activation_solution}
  \boxed{N(t)=\frac{R}{\lambda}\left(1-e^{-\lambda t}\right).}
\end{equation}
The activity during production is
\begin{equation}
  \label{eq:lec05:activation_activity}
  \boxed{\activity(t)=\lambda N(t)=R\left(1-e^{-\lambda t}\right).}
\end{equation}
\textit{Physical significance:} Production saturates.  After several lifetimes,
the source is producing and losing nuclei at the same rate, so the activity
approaches \(R\).  Irradiating much longer does not increase the activity
substantially.

%%───────────────────────────────────────────────────────────────────
\subsection{Decay Chains and Equilibrium}
\label{subsec:lec05:decay_chains}

\subsubsection*{Physical Motivation}
Often a radioactive parent decays into a daughter nucleus that is itself
radioactive.  Then the daughter abundance is controlled by two competing
effects: it is produced by parent decay and destroyed by its own decay.

Consider a parent \(1\) with decay constant \(\lambda_1\), initially
\(N_1(0)=N_0\), decaying into a daughter \(2\) with decay constant
\(\lambda_2\), initially \(N_2(0)=0\):
\begin{align}
  \label{eq:lec05:chain_parent}
  \dv{N_1}{t} &= -\lambda_1 N_1,\\
  \label{eq:lec05:chain_daughter}
  \dv{N_2}{t} &= \lambda_1N_1-\lambda_2N_2.
\end{align}
The parent solution is
\begin{equation}
  \label{eq:lec05:chain_parent_solution}
  \boxed{N_1(t)=N_0e^{-\lambda_1t}.}
\end{equation}
Substituting this into Eq.~\eqref{eq:lec05:chain_daughter} gives
\begin{equation}
  \label{eq:lec05:chain_daughter_solution}
  \boxed{N_2(t)=\frac{\lambda_1N_0}{\lambda_2-\lambda_1}
  \left(e^{-\lambda_1t}-e^{-\lambda_2t}\right)}
  \qquad(\lambda_1\neq\lambda_2).
\end{equation}
The daughter activity is therefore
\begin{equation}
  \label{eq:lec05:chain_daughter_activity}
  \boxed{\activity_2(t)=\lambda_2N_2(t)
  =\frac{\lambda_1\lambda_2N_0}{\lambda_2-\lambda_1}
  \left(e^{-\lambda_1t}-e^{-\lambda_2t}\right).}
\end{equation}

\paragraph{Long-lived parent, short-lived daughter.}
If \(\lambda_1\ll\lambda_2\), the parent changes very slowly compared with the
daughter.  After several daughter lifetimes,
\begin{equation}
  \label{eq:lec05:secular_equilibrium}
  \boxed{\activity_2(t)\approx\activity_1(t).}
\end{equation}
\textit{Physical significance:} This is secular equilibrium.  The daughter
decays as fast as it is produced, so its activity equals the parent activity
even though the daughter may have a much shorter lifetime.

\paragraph{Comparable lifetimes.}
If \(\lambda_1<\lambda_2\) but not by a huge factor, the ratio
\(\activity_2/\activity_1\) approaches a constant after the daughter builds up:
\begin{equation}
  \label{eq:lec05:transient_equilibrium}
  \boxed{\frac{\activity_2}{\activity_1}\rightarrow
  \frac{\lambda_2}{\lambda_2-\lambda_1}.}
\end{equation}
\textit{Physical significance:} This is transient equilibrium.  Both activities
eventually fall with the parent lifetime, but their ratio becomes nearly fixed.

\paragraph{Short-lived parent, long-lived daughter.}
If \(\lambda_1>\lambda_2\), the parent disappears quickly and leaves behind a
daughter population that then decays on its own longer time scale:
\begin{equation}
  \label{eq:lec05:short_parent_limit}
  \boxed{N_2(t)\approx N_0e^{-\lambda_2t}\quad
  \text{after the parent has decayed}.}
\end{equation}
\textit{Physical significance:} The source changes character: initially the
parent dominates the radiation, but later the daughter controls the activity.

%%───────────────────────────────────────────────────────────────────
\subsection{Radiation and Health}
\label{subsec:lec05:radiation_health}

\subsubsection*{Physical Motivation}
Activity tells us how many decays occur per second, but biological risk depends
on what radiation reaches tissue, how much energy it deposits, and how damaging
that radiation type is.  We therefore need dose units, not only activity units.

\paragraph{Ionizing versus non-ionizing radiation.}
Ionizing radiation can remove electrons from atoms and break molecular bonds,
including bonds in DNA.  Non-ionizing radiation does not have a comparably clear
ionization mechanism at ordinary exposure levels, although very intense
non-ionizing radiation can still cause heating and other effects.

\paragraph{Radiation-protection principles.}
Two working principles are used in radiation protection:
\begin{itemize}
  \item \textbf{ALARA}: exposures should be kept \emph{As Low As Reasonably
  Achievable}.  This is a practical rule, not a microscopic law.
  \item \textbf{LNT}: the \emph{Linear No-Threshold} model assumes that risk is
  proportional to dose even at low doses, with no perfectly safe threshold.
\end{itemize}
\textit{Physical significance:} LNT is a conservative regulatory model.  It is
used because high-dose harm is known and low-dose risk is hard to measure
directly; it should not be confused with a precisely proven microscopic law at
all dose levels.

%%───────────────────────────────────────────────────────────────────
\subsection{Radiation Units}
\label{subsec:lec05:radiation_units}

\subsubsection*{Physical Motivation}
Different questions require different units.  ``How many nuclei decay per
second?'' is an activity question.  ``How much energy entered the material?'' is
an absorbed-dose question.  ``How risky is this for tissue?'' is an equivalent
dose question.

\paragraph{Exposure.}
Exposure measures ionization produced in air.  The old unit is the roentgen:
\begin{equation}
  \label{eq:lec05:roentgen}
  \boxed{1\,\mathrm{R}=2.58\times10^{-4}\,\mathrm{C\,kg^{-1}}
  \quad\text{in air}.}
\end{equation}
It is historically important, but less fundamental than absorbed dose.

\paragraph{Absorbed dose.}
Absorbed dose \(D\) measures deposited energy per unit mass:
\begin{equation}
  \label{eq:lec05:gray}
  \boxed{1\,\mathrm{Gy}=1\,\mathrm{J\,kg^{-1}},\qquad
  1\,\mathrm{rad}=10^{-2}\,\mathrm{Gy}.}
\end{equation}
\textit{Physical significance:} Absorbed dose is a physical energy-deposition
quantity.  It does not yet say whether the radiation was alpha, beta, gamma,
neutrons, or protons.

\paragraph{Equivalent dose.}
Equivalent dose \(H\) weights the absorbed dose by a radiation quality factor
\(Q\):
\begin{equation}
  \label{eq:lec05:equivalent_dose}
  \boxed{H=Q\,D,\qquad 1\,\mathrm{Sv}=1\,\mathrm{J\,kg^{-1}},\qquad
  1\,\mathrm{rem}=10^{-2}\,\mathrm{Sv}.}
\end{equation}
Typical quality factors are
\begin{equation}
  \label{eq:lec05:quality_factors}
  \boxed{
  Q_{\gamma}\approx Q_{\beta}\approx Q_{X}\approx1,\qquad
  Q_p\approx2,\qquad
  Q_n\approx5\text{--}20,\qquad
  Q_{\alpha}\approx20.}
\end{equation}
\textit{Physical significance:} Alpha particles are far more damaging per unit
absorbed energy than gamma rays or beta particles because they deposit their
energy densely over a short range.  This is why an external alpha source may be
easy to shield, while an internal alpha emitter can be biologically serious.

%%───────────────────────────────────────────────────────────────────
\subsection*{Key Takeaways}
%%───────────────────────────────────────────────────────────────────
\begin{itemize}
  \item Radioactive decay is probabilistic: the microscopic input is a decay
  probability per unit time, \(\lambda\).
  \item Different decay modes are different ways to lower nuclear energy while
  preserving conservation laws.
  \item Activity is \(\activity=\lambda N\), measured in becquerel or curie; it
  is not itself a dose or a risk.
  \item Half-life and mean lifetime are related by
  \(T_{1/2}=(\ln2)\tau\).
  \item In decay chains, daughters may build up to secular or transient
  equilibrium with the parent.
  \item Radiation safety uses dose concepts: absorbed dose in gray and
  equivalent dose in sievert.
\end{itemize}
